\centerline{\bf \Huge Подвешивание графа за вершину}    

\begin{flushright}
        \scriptsize{Если в ящике есть испорченный помидор, испорченный надо вытащить\\
                    Вы вообще сортировку делали? Не делал никогда? Помидоры сортировал?\\
                    Чемпион UFC (29-0-0)}
\end{flushright}

\centerline{\bf \large Для разбора}

\problem{0.1}
В Махачкале началась эпидемия дубайского шоколада. Сначала дубайским шоколадом заболел один человек. Когда человек здоров, он посещает за день всех своих больных друзей. Человек, который заболевает дубайским шоколадом, болеет ровно один день и весь день лежит дома, а на следующий день у него появляется иммунитет ровно на один день. Дубайский шоколад очень заразен, поэтому при посещении любого больного человек заболевает, если у него нет иммунитета. Докажите, что эпидемия дубайского шоколада рано или поздно остановится.

\problem{0.2}
В графе степень каждой вершины равна трём, и любые две несоседние вершины имеют общего соседа. Какое максимальное количество вершин может быть в этом графе?

\problem{0.3} 
В Дагестане из каждого административного района выходит не более 4 дорог. Также известно, что между любыми двумя районами существует путь, проходящий не более чем по двум другим районам. Докажите, что в Дагестане не более 53 района.

\centerline{\bf \large Лёгкие задачи}

\problem{1}
Докажите, что количество висячих вершин в дереве не меньше, чем максимальная из степеней его вершин.

\problem{2}
Докажите, что вершины дерева можно покрасить правильным образом в 2 цвета.

\problem{3}
Докажите, что в любом связном графе найдётся вершина, при удалении которой(вместе с исходящими из нее ребрами) оставшийся граф будет связным.

\problem{4}
Степень всех вершин графа не меньше $n$, причем в нем нет циклов длины $3,4,5$. Докажите, что в нем найдутся $n^2-n$ вершин, никакие две из которых не смежны.

\problem{5}
Степень любой вершины графа не превосходит $n$. Докажите, что его можно покрасить в $n +1$   цвет правильным образом.

\problem{6}
В стране 125   городов, некоторые из которых соединены друг с другом дорогами. Из каждого города выходит хотя бы 5   дорог. Докажите, что существует несамопересекающийся циклический маршрут, состоящий из не более чем 6   городов.


\newpage

\hspace{3mm}


\centerline{\bf \large Сложные задачи}

\problem{1}
В классе 30 человек, один из них Ахмеда. Каждый из Ахмеда одноклассников имеет ровно пять общих друзей с Ахмедом. Докажите, что в классе есть ученик с нечётным числом друзей.

\problem{2}
В стране 100 городов, некоторые из которых соединены авиалиниями. Известно, что от каждого города можно долететь до любого другого (возможно, с пересадками). Докажите, что можно побывать во всех городах, совершив не более 196 перелётов.

\problem{3}
В смену ВУТС 2025 пришло 49 учеников. Известно, что если трое учеников попарно незнакомы друг с другом, то какие-то двое из них имеют общего знакомого. Докажите, что кто-то из учеников имеет хотя бы 6 знакомых.

\problem{4}
На планете 100 государств. Некоторые города соединены дорогами, причём из каждого города выходит не более 10 дорог. В каждом государстве есть город, из которого все дороги выходят только в города из этого государства. Докажите, что есть два города такие, что из одного нельзя попасть в другой не более чем за 2 пересадки.